% !TeX root = ../papers/nc-p3.tex
\section{モナドの退化集合の余次元と反射包}
\label{sec:finite-instanton-sheaf-boundary}
\begin{theorem}\label{thm:finite-instanton-sheaf-boundary}
定理\ref{thm:finite-symmetric-instanton-monad}の全窓で、$a$ の階数降下集合は高々一次元である。
従って同じモナドは常に捩れのない階数2のインスタントン層 $\mathcal E$ を与え、
その $H^1_*$ から窓が復元される。反射包 $F=\mathcal E^{**}$ は局所自由であり、
$Q=F/\mathcal E$ は零か純一次元で $H^*(Q(-2))=0$ である。
$d=\deg Q$ とすると $c_2(F)=4-d$、$d\le4$ で、$d=4$ なら $F\simeq\mathcal O^2$ である。
\end{theorem}
\begin{proof}[退化するシジジーの根基]
$0\ne t\in L\subset\bigwedge^2 U^*\otimes V^*$ を取り、$t$ を $V$ で縮約した二形式たちを考える。
$[s]\in\mathbb P(U)$ で $a(s)t=0$ となる条件は全二形式の根基が $s$ を含むことである。
このとき $t$ は $\bigwedge^2(s^\perp)\otimes V^*$ にあり、二形式の包は三次元全部である。
実際、その包が次元2以下なら全ての二形式が共通の一次因子 $h$ を持つ。
次元2以上でそのような因子があれば $h^\perp$ で $t$ を縮約して得る
$kh\otimes V^*$ 内の像は少なくとも二次元である。次元1なら一個の非零二形式を縮約して
$U^*\otimes kv$ 内に二次元以上を得る。いずれも $K$ に純テンソルの直線を与えて有限辺に反する。
この議論は共通の一次因子を持つ全ての場合に適用できる。
ゆえに $s$ が存在すれば一意で、三次元の二形式の包から正則に復元できる。

このような $[t]$ の集合 $T\subset\mathbb P(L)$ は、全ての二形式同士の外積を係数とする
二次式の共通零点集合である。実際この消滅条件では線型包がKlein二次曲面
$\mathrm{Gr}(2,U^*)\subset\mathbb P(\bigwedge^2 U^*)$ に入る。
そのような線型空間は共通一次因子を持つか、一つの三空間の外積二乗に入る。
これは相交わる二つの二平面を固定し、両方と交わる残りの二平面が
その交点を含むかその和に入るかを見ることで従う。前者は既に排除した。
従って $T$ 上には根基射 $s:T_{\rm red}\to\mathbb P(U)$ があり、像が階数降下集合である。
そのファイバーは $a(s)$ の核の射影空間である。

$T$ 全体が射影三空間、またはその中に射影平面 $D$ を含む場合を考える。
この空間上で $s^*\mathcal O(1)=\mathcal O(m)$、$H=\ker(U^*\otimes\mathcal O\to\mathcal O(m))$ とする。
全射 $V\otimes\mathcal O\to(\bigwedge^2 H)(1)$ があるので、右辺の行列式次数 $3-2m$ は非負である。
$m=0$ なら固定三行四列行列の空間内で次元3以上の線型系が全点で階数3になるが、
階数2以下の行列は余次元2なので不可能である。従って $m=1$ で、根基射は線型埋込みとなる。
像のベクトル空間を $B\subset U$ とする。$\dim B\ge3$ で、$B$ 上の線型族は
$t(s)(v)=\iota_s\omega_v$、$\omega_v\in\bigwedge^3U^*$ と書ける。
実際 $\iota_s t(s)=0$ を偏極すると一つの引数が $B$ に入る三重線型値は完全交代であり、
補空間は高々一次元なので全三形式に延びる。体積形式で $\omega_v$ を $U$ の元と読むと
$V\to U$ は同型である。像が次元2以下、または三次元で $s$ がその像と $B$ の交叉にあるとき、
その $s$ で二形式の包が次元3に達しないからである。
従って全ての $s\in B, u\in U$ の縮約が $K$ に与える空間は二文字の全交代空間となる。
実際 $B\wedge U=\bigwedge^2U$ である。隣接対も対称になり、既知の反対対称性と合わせて
三位置対称、従って $\rho=6$ となる矛盾を得る。これで $T$ は全空間でもなく平面を含まない。

二次式の基点なので残る曲面成分は既約二次曲面のみである。その成分を $D$ とし、像が曲面だと仮定する。
滑らかな場合には $s^*\mathcal O(1)=\mathcal O(m,n)$ として同じ行列式を両母線で読み $m,n\le1$。
曲面像には $(m,n)=(1,1)$ が必要である。錐の場合も超平面束となる。
実際、頂点を解消した $\mathbb F_2$ の負切断 $C_0$ とファイバー $f$ に関し、
引戻しの因子は $C_0$ 上の次数0から $m(C_0+2f)$ であり、母線上で $m\le1$ である。
よってどちらも根基射は周囲の線型写像の制限で、曲面像ならその階数は4である。
階数3では基点を避けた点から平面への二次射影で、標数零の一般点は二つの逆像を持つ。
それらを含む核の射影直線も $T$ に入るが、一般ファイバーは零次元である
（$T$ は高々二次元で、曲面像を持つ）。これが矛盾である。
この線型同定で $s$ を変数にすると $\iota_s t(s)$ は $D$ の方程式の定数ベクトル倍である。
さらに $s$ で縮約すると恒等零なのでこの定数ベクトルも零である。すると $T$ が全空間となってしまう。
従って退化集合は高々一次元である。
\end{proof}

\begin{proof}[特異商の次数]
$a$ は層として単射であり、$\mathcal E$ は局所的に自由分解長1以下である。
余次元1で局所自由、余次元2以上では深さが少なくとも1なので捩れがない。
行列式は自明なため、$F\simeq\mathcal E^*$ である（局所自由な開集合での外積対合を反射延長する）。
$Q$ の台は高々一次元。層の双対複体から $F$ は
$J=\operatorname{coker}(b^\vee)$ の部分層、商は $4\mathcal O(1)$ の部分層である。
$0\to4\mathcal O(-1)\to10\mathcal O\to J\to0$ により、
$H^0(F(-1))=0$、$H^1(F(t))=0$ ($t\le-2$) を得る。
$H^2(\mathcal E(-2))=0$ は元の二つの短完全列から従うので $H^2(F(-2))=0$ でもある。
またSerre双対により $H^3(F(-2))=0$。Riemann--Rochより $c_3(F)=2\chi(F(-2))=0$。
これは反射階数2層の点特異性がないことを意味する。実際その局所射影次元は高々1で、
双対化した自由分解のChern指標は $[F^*]-[\mathcal Ext^1(F,\mathcal O)]$ のものである。
$F^*\simeq F$ と第三指標の符号を比較すると点支持の後者の長さは $c_3(F)=0$ となる。
局所極小表示により $F$ は局所自由である。
$Q\ne0$ なら閉点で $\mathcal E$ の深さは2以上、$F$ は自由なので $Q$ の深さは1以上、
従って純一次元である。$\mathcal E(-2),F(-2)$ の全コホモロジーは零なので $Q(-2)$ も同じである。
$d=c_2(\mathcal E)-c_2(F)$ である。$F^*=F$ と既に示した消滅から
$h^1(F(-1))=c_2(F)$ なので $d\le4$。等号なら
$H^1(F(-1))=H^2(F(-2))=H^3(F(-3))=0$ より $F$ は0正則で大域生成される。
自明な行列式と階数2より $F$ は自明である。
\end{proof}
